% answers.tex —— 2014 年·辽宁卷·理科 · 参考答案与详细解析片段
% 由 tools/generate_projects.py 依来源总账自 2014/2014_数学真题_辽宁理_试题.md + 2014/2014_数学真题_辽宁理_解析正文.md + 2014/2014_数学真题_辽宁理_答案要点.md 生成，勿手改（改源 MD 后重跑）


\begin{answer}
\gksolhead{1}{}

\textbf{题目}：

已知全集 $U=\mathbf R$，$A=\{x\mid x\leqslant 0\}$，$B=\{x\mid x\geqslant 1\}$，则集合 $\complement_U(A\cup B)=$\mbox{（　　）}

（A）$\{x\mid x\geqslant 0\}$；（B）$\{x\mid x\leqslant 1\}$；（C）$\{x\mid 0\leqslant x\leqslant 1\}$；（D）$\{x\mid 0<x<1\}$。

\Needspace{4\baselineskip}
\textbf{解析}

\Needspace{5\baselineskip}
先写出两个集合：
\nopagebreak[4]
\[A=(-\infty,0],\qquad B=[1,+\infty).\]
\Needspace{5\baselineskip}
所以
\nopagebreak[4]
\[A\cup B=(-\infty,0]\cup[1,+\infty).\]
\Needspace{5\baselineskip}
在数轴上取补集，剩下的恰是 $0$ 与 $1$ 之间的部分，且两个端点都不属于 $A\cup B$，故
\nopagebreak[4]
\[\complement_U(A\cup B)=(0,1).\]

选 D。

\gkanswer{答案：D。$A\cup B=(-\infty,0]\cup[1,+\infty)$，其补集为 $(0,1)$。}

\end{answer}
\begin{answer}
\gksolhead{2}{}

\textbf{题目}：

设复数 $z$ 满足 $(z-2\mathrm i)(2-\mathrm i)=5$，则 $z=$\mbox{（　　）}

（A）$2+3\mathrm i$；（B）$2-3\mathrm i$；（C）$3+2\mathrm i$；（D）$3-2\mathrm i$。

\Needspace{4\baselineskip}
\textbf{解析}

\Needspace{5\baselineskip}
由 $(z-2\mathrm i)(2-\mathrm i)=5$ 得
\nopagebreak[4]
\[z-2\mathrm i=\frac5{2-\mathrm i}.\]
\Needspace{5\baselineskip}
分子分母同乘 $2+\mathrm i$：
\nopagebreak[4]
\[\frac5{2-\mathrm i}=\frac{5(2+\mathrm i)}{(2-\mathrm i)(2+\mathrm i)}
=\frac{5(2+\mathrm i)}{4+1}=2+\mathrm i.\]
\Needspace{5\baselineskip}
于是
\nopagebreak[4]
\[z=2+\mathrm i+2\mathrm i=2+3\mathrm i.\]

选 A。

\gkanswer{答案：A。$z-2\mathrm i=\dfrac5{2-\mathrm i}=2+\mathrm i$，故 $z=2+3\mathrm i$。}

\end{answer}
\begin{answer}
\gksolhead{3}{}

\textbf{题目}：

已知 $a=2^{-\frac13}$，$b=\log_2\dfrac13$，$c=\log_{\frac12}\dfrac13$，则\mbox{（　　）}

（A）$a>b>c$；（B）$a>c>b$；（C）$c>a>b$；（D）$c>b>a$。

\Needspace{4\baselineskip}
\textbf{解析}

三个数分别与 $0$、$1$ 比较。

\Needspace{5\baselineskip}
因为 $2^{\frac13}>2^0=1$，所以
\nopagebreak[4]
\[0<a=2^{-\frac13}<1.\]
\Needspace{5\baselineskip}
由对数性质
\nopagebreak[4]
\[b=\log_2\frac13=-\log_23<0,\]
\Needspace{5\baselineskip}
而
\nopagebreak[4]
\[c=\log_{\frac12}\frac13=\frac{\ln\frac13}{\ln\frac12}
=\frac{-\ln3}{-\ln2}=\log_23>1.\]
\Needspace{5\baselineskip}
因此
\nopagebreak[4]
\[c>1>a>0>b,\]
即 $c>a>b$。

选 C。

\gkanswer{答案：C。$0<a=2^{-\frac13}<1$，$b=\log_2\dfrac13<0$，$c=\log_{\frac12}\dfrac13=\log_23>1$，故 $c>a>b$。}

\end{answer}
\begin{answer}
\gksolhead{4}{}

\textbf{题目}：

已知 $m$，$n$ 表示两条不同直线，$\alpha$ 表示平面，下列说法正确的是\mbox{（　　）}

（A）若 $m\parallel\alpha$，$n\parallel\alpha$，则 $m\parallel n$；（B）若 $m\perp\alpha$，$n\subset\alpha$，则 $m\perp n$；（C）若 $m\perp\alpha$，$m\perp n$，则 $n\parallel\alpha$；（D）若 $m\parallel\alpha$，$m\perp n$，则 $n\perp\alpha$。

\Needspace{4\baselineskip}
\textbf{解析}

逐项判断。

（A）$m\parallel\alpha$ 且 $n\parallel\alpha$ 时，$m$ 与 $n$ 可能相交，也可能异面，不一定平行，故（A）错。

（B）若 $m\perp\alpha$，则 $m$ 垂直于平面 $\alpha$ 内的每一条直线；由 $n\subset\alpha$ 得 $m\perp n$，故（B）正确。

（C）若 $m\perp\alpha$，$m\perp n$，则 $n$ 可能平行于 $\alpha$，也可能在 $\alpha$ 内，故（C）错。

（D）取 $\alpha$ 为 $xOy$ 坐标平面，$m$ 为 $x$ 轴所在直线，$n$ 为 $y$ 轴所在直线，则 $m\parallel\alpha$，$m\perp n$，但 $n\subset\alpha$，$n$ 不垂直于 $\alpha$，故（D）错。

选 B。

\gkanswer{答案：B。线面垂直则垂直于面内任一直线；其余三项均可由反例否定。}

\end{answer}
\begin{answer}
\gksolhead{5}{}

\textbf{题目}：

设 $\boldsymbol a$，$\boldsymbol b$，$\boldsymbol c$ 是非零向量，已知命题 $p$：若 $\boldsymbol a\cdot\boldsymbol b=0$，$\boldsymbol b\cdot\boldsymbol c=0$，则 $\boldsymbol a\cdot\boldsymbol c=0$；命题 $q$：若 $\boldsymbol a\parallel\boldsymbol b$，$\boldsymbol b\parallel\boldsymbol c$，则 $\boldsymbol a\parallel\boldsymbol c$，则下列命题中的真命题是\mbox{（　　）}

（A）$p\vee q$；（B）$p\wedge q$；（C）$(\neg p)\wedge(\neg q)$；（D）$p\vee(\neg q)$。

\Needspace{4\baselineskip}
\textbf{解析}

先判断 $p$、$q$ 的真假。

\Needspace{5\baselineskip}
对 $p$，取
\nopagebreak[4]
\[\boldsymbol a=(1,0,0),\qquad \boldsymbol b=(0,1,0),\qquad \boldsymbol c=(1,0,1),\]
\Needspace{5\baselineskip}
则
\nopagebreak[4]
\[\boldsymbol a\cdot\boldsymbol b=0,\qquad \boldsymbol b\cdot\boldsymbol c=0,\]
\Needspace{5\baselineskip}
但
\nopagebreak[4]
\[\boldsymbol a\cdot\boldsymbol c=1\neq0.\]
所以 $p$ 是假命题。

\Needspace{5\baselineskip}
对 $q$，因为 $\boldsymbol b\neq\boldsymbol 0$，由 $\boldsymbol a\parallel\boldsymbol b$ 可设 $\boldsymbol a=\lambda\boldsymbol b$，由 $\boldsymbol b\parallel\boldsymbol c$ 可设 $\boldsymbol c=\mu\boldsymbol b$，其中 $\lambda,\mu\neq0$。于是
\nopagebreak[4]
\[\boldsymbol a=\frac{\lambda}{\mu}\boldsymbol c,\]
故 $\boldsymbol a\parallel\boldsymbol c$，$q$ 是真命题。

于是 $p\vee q$ 真，$p\wedge q$ 假，$(\neg p)\wedge(\neg q)$ 假，$p\vee(\neg q)$ 假。

选 A。

\gkanswer{答案：A。取 $\boldsymbol a=(1,0,0)$，$\boldsymbol b=(0,1,0)$，$\boldsymbol c=(1,0,1)$ 知 $p$ 假；$q$ 真，故 $p\vee q$ 真。}

\end{answer}
\begin{answer}
\gksolhead{6}{}

\textbf{题目}：

6 把椅子摆成一排，3 人随机就座，任何两人不相邻的坐法种数为\mbox{（　　）}

（A）$144$；（B）$120$；（C）$72$；（D）$24$。

\Needspace{4\baselineskip}
\textbf{解析}

先从 6 个座位中选出 3 个两两不相邻的座位。

\Needspace{5\baselineskip}
把选中的座位记为 $x_1<x_2<x_3$。两两不相邻等价于
\nopagebreak[4]
\[x_2\geqslant x_1+2,\qquad x_3\geqslant x_2+2.\]
\Needspace{5\baselineskip}
令
\nopagebreak[4]
\[y_1=x_1,\qquad y_2=x_2-1,\qquad y_3=x_3-2,\]
\Needspace{5\baselineskip}
则 $1\leqslant y_1<y_2<y_3\leqslant 4$，即从 $1,2,3,4$ 中选 3 个，共有
\nopagebreak[4]
\[\mathrm C_4^3=4\]
种选法（等价于插空法：先排 3 个空位，形成 4 个空，选 3 个放人）。

\Needspace{5\baselineskip}
选定座位后，3 个人在 3 个座位上的排列有
\nopagebreak[4]
\[3!=6\]
\Needspace{5\baselineskip}
种。故所求坐法种数为
\nopagebreak[4]
\[4\times6=24.\]

选 D。

\gkanswer{答案：D。先从 6 个座位中选 3 个两两不相邻的座位，有 $\mathrm C_4^3=4$ 种，再排 3 人，有 $3!=6$ 种，共 $4\times6=24$ 种。}

\end{answer}
\begin{answer}
\gksolhead{7}{}

\textbf{题目}：

某几何体三视图如图所示，则该几何体的体积为\mbox{（　　）}
\nopagebreak[4]
\begin{center}% ============================================================
% f2014_liaoning_science_q07_three_views.tex
% 2014 辽宁卷（理科）第 7 题：三视图
% 来源 PDF：普通高考/2014/2014辽宁理.pdf
% 原页：第 1 页（A3），印面 895，第 7 题题干下方
% 主视图、左视图均为 2×2 矩形，中部有竖线；
% 俯视图为 2×2 正方形去掉左上、右下两个半径 1 的四分之一圆。
% ============================================================
\begin{tikzpicture}[>=Stealth, line width=0.7pt, font=\small,
  dim/.style={line width=0.4pt},
  every node/.style={inner sep=1.5pt}]
  % ---------- 主视图 ----------
  \draw (0,3.0) rectangle (2,5.0);
  \draw (1,3.0) -- (1,5.0);
  \node at (1,2.62) {主视图};
  \draw[dim,<->] (0,5.2) -- (1,5.2) node[midway,above,font=\scriptsize] {$1$};
  \draw[dim,<->] (0,5.6) -- (2,5.6) node[midway,above,font=\scriptsize] {$2$};
  \draw[dim,<->] (-0.3,5.0) -- (-0.3,3.0) node[midway,left,font=\scriptsize] {$2$};

  % ---------- 左视图 ----------
  \draw (3.6,3.0) rectangle (5.6,5.0);
  \draw (4.6,3.0) -- (4.6,5.0);
  \node at (4.6,2.62) {左视图};
  \draw[dim,<->] (3.6,5.2) -- (4.6,5.2) node[midway,above,font=\scriptsize] {$1$};
  \draw[dim,<->] (3.6,5.6) -- (5.6,5.6) node[midway,above,font=\scriptsize] {$2$};

  % ---------- 俯视图 ----------
  % 正方形去掉左上、右下两个四分之一圆（半径 1）
  \draw (0,0) -- (1,0)
        arc[start angle=180, end angle=90, radius=1]
        -- (2,2) -- (1,2)
        arc[start angle=0, end angle=-90, radius=1]
        -- cycle;
  \node at (1,-0.95) {俯视图};
  \draw[dim,<->] (0,-0.32) -- (2,-0.32) node[midway,below,font=\scriptsize] {$2$};
  \draw[dim,<->] (2.35,0) -- (2.35,2) node[midway,right,font=\scriptsize] {$2$};
  % 两个半径标注
  \draw[dim,->] (0,2) -- (0.707,1.293) node[midway,above left,font=\scriptsize] {$1$};
  \draw[dim,->] (2,0) -- (1.293,0.707) node[midway,below right,font=\scriptsize] {$1$};
\end{tikzpicture}
\end{center}
\nopagebreak[4]
（A）$8-2\pi$；（B）$8-\pi$；（C）$8-\dfrac\pi2$；（D）$8-\dfrac\pi4$。

\Needspace{4\baselineskip}
\textbf{解析}

由三视图还原几何体。

主视图、左视图都是边长 $2$ 的矩形，说明几何体是高为 $2$ 的柱体，底面就是俯视图所画的图形。

\Needspace{5\baselineskip}
俯视图是一个边长为 $2$ 的正方形，在左上、右下两个角上各去掉一个半径为 $1$ 的四分之一圆。其面积为
\nopagebreak[4]
\[S=2\times2-2\times\frac14\pi\cdot1^2=4-\frac\pi2.\]
\Needspace{5\baselineskip}
柱体高 $h=2$，故体积
\nopagebreak[4]
\[V=Sh=\left(4-\frac\pi2\right)\times2=8-\pi.\]

选 B。

\gkanswer{答案：B。俯视图面积 $4-\dfrac\pi2$，柱高 $2$，故 $V=2\left(4-\dfrac\pi2\right)=8-\pi$。}

\end{answer}
\begin{answer}
\gksolhead{8}{}

\textbf{题目}：

设等差数列 $\{a_n\}$ 的公差为 $d$，若数列 $\{2^{a_1a_n}\}$ 为递减数列，则\mbox{（　　）}

（A）$d<0$；（B）$d>0$；（C）$a_1d>0$；（D）$a_1d<0$。

\Needspace{4\baselineskip}
\textbf{解析}

\Needspace{5\baselineskip}
记 $b_n=2^{a_1a_n}$，则
\nopagebreak[4]
\[\frac{b_{n+1}}{b_n}
=2^{a_1a_{n+1}-a_1a_n}
=2^{a_1(a_{n+1}-a_n)}
=2^{a_1d}.\]
\Needspace{5\baselineskip}
数列 $\{b_n\}$ 递减等价于对一切 $n$ 有
\nopagebreak[4]
\[\frac{b_{n+1}}{b_n}<1,\]
\Needspace{5\baselineskip}
即
\nopagebreak[4]
\[2^{a_1d}<1=2^0.\]
\Needspace{5\baselineskip}
由指数函数单调递增得
\nopagebreak[4]
\[a_1d<0.\]

选 D。

\gkanswer{答案：D。相邻项之比 $2^{a_1d}<1$，故 $a_1d<0$。}

\end{answer}
\begin{answer}
\gksolhead{9}{}

\textbf{题目}：

将函数 $y=3\sin\left(2x+\dfrac\pi3\right)$ 的图象向右平移 $\dfrac\pi2$ 个单位长度，所得图象对应的函数\mbox{（　　）}

（A）在区间 $\left[\dfrac\pi{12},\dfrac{7\pi}{12}\right]$ 上单调递减；（B）在区间 $\left[\dfrac\pi{12},\dfrac{7\pi}{12}\right]$ 上单调递增；（C）在区间 $\left[-\dfrac\pi6,\dfrac\pi3\right]$ 上单调递减；（D）在区间 $\left[-\dfrac\pi6,\dfrac\pi3\right]$ 上单调递增。

\Needspace{4\baselineskip}
\textbf{解析}

\Needspace{5\baselineskip}
设平移后所得函数为 $h(x)$，则
\nopagebreak[4]
\[\begin{aligned}
h(x)&=3\sin\left[2\left(x-\frac\pi2\right)+\frac\pi3\right]\\
&=3\sin\left(2x-\pi+\frac\pi3\right)\\
&=3\sin\left(2x-\frac{2\pi}3\right).
\end{aligned}\]
\Needspace{5\baselineskip}
令
\nopagebreak[4]
\[-\frac\pi2+2k\pi\leqslant 2x-\frac{2\pi}3\leqslant\frac\pi2+2k\pi,\]
\Needspace{5\baselineskip}
解得
\nopagebreak[4]
\[\frac\pi{12}+k\pi\leqslant x\leqslant\frac{7\pi}{12}+k\pi,\quad k\in\mathbf Z.\]
\Needspace{5\baselineskip}
取 $k=0$，得 $h(x)$ 的一个单调递增区间为
\nopagebreak[4]
\[\left[\frac\pi{12},\frac{7\pi}{12}\right].\]
故（B）正确，（A）错误。

\Needspace{5\baselineskip}
当 $x\in\left[-\dfrac\pi6,\dfrac\pi3\right]$ 时，
\nopagebreak[4]
\[2x-\frac{2\pi}3\in\left[-\pi,0\right].\]
函数 $y=3\sin u$ 在 $\left[-\pi,-\dfrac\pi2\right]$ 上递减、在 $\left[-\dfrac\pi2,0\right]$ 上递增，故 $h(x)$ 在该区间上不单调，（C）（D）均错误。

选 B。

\gkanswer{答案：B。平移后 $y=3\sin\left(2x-\dfrac{2\pi}3\right)$，在 $\left[\dfrac\pi{12},\dfrac{7\pi}{12}\right]$ 上单调递增。}

\end{answer}
\begin{answer}
\gksolhead{10}{}

\textbf{题目}：

已知点 $A(-2,3)$ 在抛物线 $C:y^2=2px$ 的准线上，过点 $A$ 的直线与 $C$ 在第一象限相切于点 $B$，记 $C$ 的焦点为 $F$，则直线 $BF$ 的斜率为\mbox{（　　）}

（A）$\dfrac12$；（B）$\dfrac23$；（C）$\dfrac34$；（D）$\dfrac43$。

\Needspace{4\baselineskip}
\textbf{解析}

\Needspace{5\baselineskip}
抛物线 $y^2=2px$ 的准线为
\nopagebreak[4]
\[x=-\frac p2.\]
\Needspace{5\baselineskip}
点 $A(-2,3)$ 在准线上，故
\nopagebreak[4]
\[-\frac p2=-2,\qquad p=4.\]
\Needspace{5\baselineskip}
抛物线为
\nopagebreak[4]
\[C:y^2=8x,\qquad F(2,0).\]

\Needspace{5\baselineskip}
设切点为 $B(x_0,y_0)$，$y_0>0$。对 $y^2=8x$ 两边求导得
\nopagebreak[4]
\[2y\,y'=8,\qquad y'=\frac4y,\]
\Needspace{5\baselineskip}
故切线斜率为 $\dfrac4{y_0}$，切线方程为
\nopagebreak[4]
\[y-y_0=\frac4{y_0}(x-x_0).\]
\Needspace{5\baselineskip}
它过 $A(-2,3)$，代入得
\nopagebreak[4]
\[3-y_0=\frac4{y_0}(-2-x_0).\]
\Needspace{5\baselineskip}
又 $y_0^2=8x_0$，即 $x_0=\dfrac{y_0^2}{8}$。代入整理：
\nopagebreak[4]
\[\begin{aligned}
3-y_0&=\frac4{y_0}\left(-2-\frac{y_0^2}{8}\right)
=-\frac8{y_0}-\frac{y_0}{2},\\
6y_0-2y_0^2&=-16-y_0^2,\\
y_0^2-6y_0-16&=0.
\end{aligned}\]
\Needspace{5\baselineskip}
解得 $y_0=8$ 或 $y_0=-2$。因为 $B$ 在第一象限，取 $y_0=8$，于是
\nopagebreak[4]
\[x_0=\frac{64}{8}=8,\qquad B(8,8).\]
\Needspace{5\baselineskip}
所以
\nopagebreak[4]
\[k_{BF}=\frac{8-0}{8-2}=\frac43.\]

选 D。

\gkanswer{答案：D。准线 $x=-\dfrac p2=-2$ 得 $p=4$，切点 $B(8,8)$，$F(2,0)$，故 $k_{BF}=\dfrac{8}{6}=\dfrac43$。}

\end{answer}
\begin{answer}
\gksolhead{11}{}

\textbf{题目}：

当 $x\in[-2,1]$ 时，不等式 $ax^3-x^2+4x+3\geqslant 0$ 恒成立，则实数 $a$ 的取值范围是\mbox{（　　）}

（A）$[-5,-3]$；（B）$\left[-6,-\dfrac98\right]$；（C）$[-6,-2]$；（D）$[-4,-3]$。

\Needspace{4\baselineskip}
\textbf{解析}

当 $x=0$ 时，不等式为 $3\geqslant0$，恒成立。

\Needspace{5\baselineskip}
当 $x\in(0,1]$ 时，$x^3>0$，原不等式等价于
\nopagebreak[4]
\[a\geqslant\frac{x^2-4x-3}{x^3}.\]
\Needspace{5\baselineskip}
记
\nopagebreak[4]
\[g(x)=\frac{x^2-4x-3}{x^3}=\frac1x-\frac4{x^2}-\frac3{x^3},\]
\Needspace{5\baselineskip}
则
\nopagebreak[4]
\[g'(x)=-\frac1{x^2}+\frac8{x^3}+\frac9{x^4}
=\frac{9+8x-x^2}{x^4}>0\quad(0<x\leqslant1),\]
\Needspace{5\baselineskip}
故 $g$ 在 $(0,1]$ 上单调递增，
\nopagebreak[4]
\[g(x)\leqslant g(1)=-6,\]
\Needspace{5\baselineskip}
所以
\nopagebreak[4]
\[a\geqslant-6.\]

\Needspace{5\baselineskip}
当 $x\in[-2,0)$ 时，$x^3<0$，原不等式等价于
\nopagebreak[4]
\[a\leqslant\frac{x^2-4x-3}{x^3}=g(x).\]
\Needspace{5\baselineskip}
此时
\nopagebreak[4]
\[g'(x)=\frac{9+8x-x^2}{x^4}=-\frac{(x-9)(x+1)}{x^4}.\]
\Needspace{5\baselineskip}
在 $[-2,-1)$ 上 $(x-9)<0$，$(x+1)<0$，故 $g'(x)<0$，$g$ 递减；在 $(-1,0)$ 上 $g'(x)>0$，$g$ 递增。所以 $g$ 在 $[-2,0)$ 上的最小值为
\nopagebreak[4]
\[g(-1)=\frac{1+4-3}{-1}=-2,\]
\Needspace{5\baselineskip}
从而
\nopagebreak[4]
\[a\leqslant-2.\]

\Needspace{5\baselineskip}
综上
\nopagebreak[4]
\[-6\leqslant a\leqslant-2.\]

选 C。

\gkanswer{答案：C。分离参数得 $a\geqslant-6$（$x>0$）且 $a\leqslant-2$（$x<0$），故 $a\in[-6,-2]$。}

\end{answer}
\begin{answer}
\gksolhead{12}{}

\textbf{题目}：

已知定义在 $[0,1]$ 上的函数 $f(x)$ 满足：① $f(0)=f(1)=0$；② 对所有 $x$，$y\in[0,1]$，且 $x\neq y$，有 $|f(x)-f(y)|<\dfrac12|x-y|$。若对所有 $x$，$y\in[0,1]$，$|f(x)-f(y)|<k$ 恒成立，则 $k$ 的最小值为\mbox{（　　）}

（A）$\dfrac12$；（B）$\dfrac14$；（C）$\dfrac1{2\pi}$；（D）$\dfrac18$。

\Needspace{4\baselineskip}
\textbf{解析}

\Needspace{5\baselineskip}
对任意 $x\in[0,1]$，由 $f(0)=0$ 及条件②，
\nopagebreak[4]
\[|f(x)|=|f(x)-f(0)|<\frac12|x-0|=\frac x2.\]
\Needspace{5\baselineskip}
由 $f(1)=0$ 同理，
\nopagebreak[4]
\[|f(x)|=|f(x)-f(1)|<\frac12|1-x|=\frac{1-x}2.\]
\Needspace{5\baselineskip}
因此
\nopagebreak[4]
\[|f(x)|<\frac12\min\{x,1-x\}\leqslant\frac14.\]

\Needspace{5\baselineskip}
再估计两点之差。设 $0\leqslant x<y\leqslant1$。一方面，
\nopagebreak[4]
\[|f(y)-f(x)|<\frac12(y-x).\]
\Needspace{5\baselineskip}
另一方面，由 $f(0)=f(1)=0$，
\nopagebreak[4]
\[\begin{aligned}
f(y)-f(x)
&=\bigl(f(y)-f(1)\bigr)-\bigl(f(x)-f(0)\bigr)\\
&<\frac12(1-y)+\frac12x
=\frac12(1-y+x).
\end{aligned}\]
\Needspace{5\baselineskip}
由于
\nopagebreak[4]
\[(y-x)+(1-y+x)=1,\]
\Needspace{5\baselineskip}
所以
\nopagebreak[4]
\[\min\{y-x,\ 1-y+x\}\leqslant\frac12,\]
\Needspace{5\baselineskip}
从而
\nopagebreak[4]
\[|f(x)-f(y)|<\frac12\cdot\frac12=\frac14.\]
故 $k=\dfrac14$ 一定满足要求。

\Needspace{5\baselineskip}
再说明不能更小。取
\nopagebreak[4]
\[f_\varepsilon(x)=\left(\frac12-\varepsilon\right)\min\{x,1-x\},\qquad 0<\varepsilon<\frac12.\]
\Needspace{5\baselineskip}
它满足 $f_\varepsilon(0)=f_\varepsilon(1)=0$，且斜率绝对值不超过 $\dfrac12-\varepsilon<\dfrac12$，故满足条件②。但
\nopagebreak[4]
\[\left|f_\varepsilon\left(\frac12\right)-f_\varepsilon(0)\right|
=\frac12\left(\frac12-\varepsilon\right)
=\frac14-\frac\varepsilon2.\]
当 $\varepsilon\to0^+$ 时该值趋于 $\dfrac14$，所以任何 $k<\dfrac14$ 都不能使 $|f(x)-f(y)|<k$ 恒成立。

因此 $k$ 的最小值为 $\dfrac14$。

选 B。

\gkanswer{答案：B。由 $f(0)=f(1)=0$ 及斜率小于 $\dfrac12$ 得 $|f(x)|<\dfrac12\min\{x,1-x\}\leqslant\dfrac14$，且上界可逼近，故 $k_{\min}=\dfrac14$。}

\end{answer}
\begin{answer}
\gksolhead{13}{}

\textbf{题目}：

执行如图的程序框图，若输入 $x=9$，则输出 $y=$\gkblank{}。
\nopagebreak[4]
\begin{center}% ============================================================
% f2014_liaoning_science_q13_flowchart.tex
% 2014 辽宁卷（理科）第 13 题：程序框图
% 来源 PDF：普通高考/2014/2014辽宁理.pdf
% 原页：第 1 页（A3），印面 895，第 13 题题干下方
% 按原图重绘：输入 x；计算 y=x/3+2；判断 |y-x|<1，
% 成立则输出 y，不成立则令 x=y 后返回重新计算 y。
% ============================================================
\begin{tikzpicture}[>=Stealth, line width=0.7pt, font=\small,
  every node/.style={inner sep=2pt},
  term/.style={draw, rounded corners=3pt, minimum width=2.1cm,
    minimum height=0.6cm, align=center},
  proc/.style={draw, minimum width=2.5cm, minimum height=0.7cm, align=center},
  io/.style={draw, trapezium, trapezium left angle=70, trapezium right angle=110,
    minimum width=2.5cm, minimum height=0.6cm, align=center},
  dec/.style={draw, diamond, aspect=1.8, minimum width=2.7cm,
    minimum height=1.0cm, align=center}]
  \node[term] (start) at (0,0) {开始};
  \node[io] (input) at (0,-0.95) {输入 $x$};
  \node[proc] (calc) at (0,-1.95) {$y=\dfrac{x}{3}+2$};
  \node[dec] (test) at (0,-3.25) {$|y-x|<1$};
  \node[io] (out) at (0,-4.55) {输出 $y$};
  \node[term] (end) at (0,-5.5) {结束};
  \node[proc] (assign) at (3.2,-1.95) {$x=y$};

  \draw[->] (start) -- (input);
  \draw[->] (input) -- (calc);
  \draw[->] (calc) -- (test);
  \draw[->] (test.south) -- node[left] {是} (out.north);
  \draw[->] (out) -- (end);
  \draw[->] (test.east) -| node[pos=0.25, above] {否} (assign.south);
  \draw[->] (assign.north) |- (0,-1.45);
\end{tikzpicture}
\end{center}
\nopagebreak[4]
\Needspace{4\baselineskip}
\textbf{解析}

按框图逐步执行。

\Needspace{5\baselineskip}
第一次：输入 $x=9$，
\nopagebreak[4]
\[y=\frac93+2=5,\qquad |y-x|=|5-9|=4\geqslant1,\]
不满足判断条件，令 $x=y=5$。

\Needspace{5\baselineskip}
第二次：
\nopagebreak[4]
\[y=\frac53+2=\frac{11}3,\qquad
\left|y-x\right|=\left|\frac{11}3-5\right|=\frac43\geqslant1,\]
仍不满足，令 $x=\dfrac{11}3$。

\Needspace{5\baselineskip}
第三次：
\nopagebreak[4]
\[y=\frac13\cdot\frac{11}3+2=\frac{11}9+2=\frac{29}9,\]
\Needspace{5\baselineskip}
此时
\nopagebreak[4]
\[\left|y-x\right|
=\left|\frac{29}9-\frac{11}3\right|
=\left|\frac{29-33}{9}\right|
=\frac49<1.\]
\Needspace{5\baselineskip}
满足判断条件，输出
\nopagebreak[4]
\[y=\frac{29}9.\]

\gkanswer{答案：$\dfrac{29}9$。依次取 $x=9,5,\dfrac{11}3$，第三次得 $y=\dfrac{29}9$，此时 $|y-x|=\dfrac49<1$，输出 $\dfrac{29}9$。}

\end{answer}
\begin{answer}
\gksolhead{14}{}

\textbf{题目}：

正方形的四个顶点 $A(-1,-1)$，$B(1,-1)$，$C(1,1)$，$D(-1,1)$ 分别在抛物线 $y=-x^2$ 和 $y=x^2$ 上，如图所示，若将一个质点随机投入正方形 $ABCD$ 中，则质点落在阴影区域的概率是\gkblank{}。
\nopagebreak[4]
\begin{center}% ============================================================
% f2014_liaoning_science_q14_parabola_region.tex
% 2014 辽宁卷（理科）第 14 题：正方形与两条抛物线
% 来源 PDF：普通高考/2014/2014辽宁理.pdf
% 原页：第 1 页（A3），印面 895，第 14 题题干下方
% 正方形 ABCD 的四个顶点分别在 y=x^2 与 y=-x^2 上，
% 阴影区域为两抛物线之间的部分。
% ============================================================
\begin{tikzpicture}[>=Stealth, line width=0.7pt, font=\small,
  x=1.9cm, y=1.9cm, every node/.style={inner sep=1.5pt}]
  % 阴影区域：-x^2 <= y <= x^2
  \fill[gray!25] (-1,1) parabola bend (0,0) (1,1)
    -- (1,-1) parabola bend (0,0) (-1,-1) -- cycle;
  % 正方形
  \draw (-1,-1) rectangle (1,1);
  % 坐标轴
  \draw[->] (-1.30,0) -- (1.45,0) node[right] {$x$};
  \draw[->] (0,-1.30) -- (0,1.45) node[above] {$y$};
  \node[below left] at (0,0) {$O$};
  % 抛物线
  \draw[line width=0.8pt, domain=-1.18:1.18, samples=80]
    plot (\x,{\x*\x});
  \draw[line width=0.8pt, domain=-1.18:1.18, samples=80]
    plot (\x,{-\x*\x});
  % 刻度
  \draw (1,0.03) -- (1,-0.03) node[below] {$1$};
  \draw (-1,0.03) -- (-1,-0.03) node[below left] {$-1$};
  \draw (0.03,1) -- (-0.03,1) node[left] {$1$};
  \draw (0.03,-1) -- (-0.03,-1) node[left] {$-1$};
  % 顶点
  \fill (-1,1) circle (1.1pt);
  \node[left] at (-1.10,1.06) {$D$};
  \fill (1,1) circle (1.1pt);
  \node[right] at (1.10,1.06) {$C$};
  \fill (-1,-1) circle (1.1pt);
  \node[left] at (-1.10,-1.06) {$A$};
  \fill (1,-1) circle (1.1pt);
  \node[right] at (1.10,-1.06) {$B$};
  % 抛物线标签
  \node[right] at (1.40,1.30) {$y=x^2$};
  \node[right] at (1.40,-1.30) {$y=-x^2$};
\end{tikzpicture}
\end{center}
\nopagebreak[4]
\Needspace{4\baselineskip}
\textbf{解析}

\Needspace{5\baselineskip}
正方形边长为 $2$，面积为
\nopagebreak[4]
\[S_{ABCD}=2^2=4.\]

\Needspace{5\baselineskip}
阴影区域是两条抛物线之间的部分，即对每个 $x\in[-1,1]$，有
\nopagebreak[4]
\[-x^2\leqslant y\leqslant x^2.\]
\Needspace{5\baselineskip}
其面积为
\nopagebreak[4]
\[\begin{aligned}
S_{\text{\text{阴影}}}
&=\int_{-1}^{1}\bigl[x^2-(-x^2)\bigr]\,\mathrm dx\\
&=\int_{-1}^{1}2x^2\,\mathrm dx
=\left.\frac{2x^3}{3}\right|_{-1}^{1}
=\frac23-\left(-\frac23\right)
=\frac43.
\end{aligned}\]
\Needspace{5\baselineskip}
这是几何概型，所求概率为
\nopagebreak[4]
\[P=\frac{S_{\text{\text{阴影}}}}{S_{ABCD}}
=\frac{4/3}{4}=\frac13.\]

\gkanswer{答案：$\dfrac13$。阴影面积 $\displaystyle\int_{-1}^{1}2x^2\,\mathrm dx=\dfrac43$，正方形面积 $4$，概率为 $\dfrac13$。}

\end{answer}
\begin{answer}
\gksolhead{15}{}

\textbf{题目}：

已知椭圆 $C:\dfrac{x^2}9+\dfrac{y^2}4=1$，点 $M$ 与 $C$ 的焦点不重合，若 $M$ 关于 $C$ 的焦点的对称点分别为 $A$，$B$，线段 $MN$ 的中点在 $C$ 上，则 $|AN|+|BN|=$\gkblank{}。

\Needspace{4\baselineskip}
\textbf{解析}

\Needspace{5\baselineskip}
椭圆中 $a=3$，$b=2$，
\nopagebreak[4]
\[c=\sqrt{a^2-b^2}=\sqrt{9-4}=\sqrt5,\]
\Needspace{5\baselineskip}
焦点为
\nopagebreak[4]
\[F_1(-\sqrt5,0),\qquad F_2(\sqrt5,0).\]

\Needspace{5\baselineskip}
因为 $A$ 是 $M$ 关于 $F_1$ 的对称点，所以 $F_1$ 是 $MA$ 的中点：
\nopagebreak[4]
\[A=2F_1-M.\]
\Needspace{5\baselineskip}
同理
\nopagebreak[4]
\[B=2F_2-M.\]
\Needspace{5\baselineskip}
设 $Q$ 为线段 $MN$ 的中点，则
\nopagebreak[4]
\[M+N=2Q.\]
\Needspace{5\baselineskip}
于是
\nopagebreak[4]
\[\begin{aligned}
|AN|
&=|N-A|
=|N-(2F_1-M)|\\
&=|M+N-2F_1|
=|2Q-2F_1|
=2|QF_1|.
\end{aligned}\]
\Needspace{5\baselineskip}
同理
\nopagebreak[4]
\[|BN|=2|QF_2|.\]
\Needspace{5\baselineskip}
因为 $Q$ 在椭圆上，由椭圆定义
\nopagebreak[4]
\[|QF_1|+|QF_2|=2a=6,\]
\Needspace{5\baselineskip}
所以
\nopagebreak[4]
\[|AN|+|BN|
=2\bigl(|QF_1|+|QF_2|\bigr)
=2\times6=12.\]

\gkanswer{答案：$12$。设 $Q$ 为 $MN$ 中点，则 $|AN|=2|F_1Q|$，$|BN|=2|F_2Q|$，故和为 $2\times2a=12$。}

\end{answer}
\begin{answer}
\gksolhead{16}{}

\textbf{题目}：

对于 $c>0$，当非零实数 $a$，$b$ 满足 $4a^2-2ab+4b^2-c=0$ 且使 $|2a+b|$ 最大时，$\dfrac3a-\dfrac4b+\dfrac5c$ 的最小值为\gkblank{}。

\Needspace{4\baselineskip}
\textbf{解析}

\Needspace{5\baselineskip}
记 $u=2a+b$。对条件左端配方：
\nopagebreak[4]
\[\begin{aligned}
4a^2-2ab+4b^2
&=\frac58(2a+b)^2+\frac38(3b-2a)^2.
\end{aligned}\]
\Needspace{5\baselineskip}
验证：右边展开为
\nopagebreak[4]
\[\frac58(4a^2+4ab+b^2)+\frac38(9b^2-12ab+4a^2)
=4a^2-2ab+4b^2.\]
\Needspace{5\baselineskip}
因为 $\dfrac38(3b-2a)^2\geqslant0$，所以
\nopagebreak[4]
\[c\geqslant\frac58(2a+b)^2,\]
\Needspace{5\baselineskip}
即
\nopagebreak[4]
\[|2a+b|\leqslant\sqrt{\frac{8c}{5}}.\]
\Needspace{5\baselineskip}
等号成立当且仅当
\nopagebreak[4]
\[3b=2a,\qquad\text{\text{即}}\quad b=\frac23a.\]
\Needspace{5\baselineskip}
故 $|2a+b|$ 最大时 $b=\dfrac23a$，此时
\nopagebreak[4]
\[c=4a^2-2a\cdot\frac23a+4\cdot\frac49a^2
=\frac{40}9a^2,\]
\Needspace{5\baselineskip}
从而
\nopagebreak[4]
\[a^2=\frac{9c}{40}.\]

\Needspace{5\baselineskip}
代入所求式子：
\nopagebreak[4]
\[\begin{aligned}
\frac3a-\frac4b+\frac5c
&=\frac3a-\frac4{2a/3}+\frac5c\\
&=\frac3a-\frac6a+\frac5c
=-\frac3a+\frac5c.
\end{aligned}\]
\Needspace{5\baselineskip}
由 $a^2=\dfrac{9c}{40}$ 得
\nopagebreak[4]
\[\frac3a=\pm\frac{3}{\sqrt{9c/40}}
=\pm\frac{2\sqrt{10}}{\sqrt c}.\]
\Needspace{5\baselineskip}
为使原式最小，取 $a>0$，即
\nopagebreak[4]
\[\frac3a=\frac{2\sqrt{10}}{\sqrt c}.\]
\Needspace{5\baselineskip}
令 $s=\dfrac1{\sqrt c}>0$，则
\nopagebreak[4]
\[\begin{aligned}
\frac3a-\frac4b+\frac5c
&=-2\sqrt{10}\,s+5s^2\\
&=5\left(s-\frac{\sqrt{10}}5\right)^2-2
\geqslant-2.
\end{aligned}\]
\Needspace{5\baselineskip}
等号在
\nopagebreak[4]
\[s=\frac{\sqrt{10}}5,\qquad c=\frac52\]
\Needspace{5\baselineskip}
时成立，此时
\nopagebreak[4]
\[a=\frac34,\qquad b=\frac12,\]
\Needspace{5\baselineskip}
且
\nopagebreak[4]
\[4a^2-2ab+4b^2=\frac94-\frac34+1=\frac52=c,\]
\nopagebreak[4]
\[|2a+b|=2=\sqrt{\frac{8c}{5}},\]
确为最大值情形。

故最小值为 $-2$。

\gkanswer{答案：$-2$。由 $4a^2-2ab+4b^2\geqslant\dfrac58(2a+b)^2$ 知取最大值时 $b=\dfrac23a$、$c=\dfrac{40}9a^2$，原式化为 $\dfrac5c-\dfrac{2\sqrt{10}}{\sqrt c}$，当 $c=\dfrac52$ 时取最小值 $-2$。}

\end{answer}
\begin{answer}
\gksolhead{17}{}

\textbf{题目}：

在 $\triangle ABC$ 中，内角 $A$，$B$，$C$ 的对边分别为 $a$，$b$，$c$，且 $a>c$，已知 $\overrightarrow{BA}\cdot\overrightarrow{BC}=2$，$\cos B=\dfrac13$，$b=3$，求：

（1） $a$ 和 $c$ 的值；

（2） $\cos(B-C)$ 的值。

\Needspace{4\baselineskip}
\textbf{解析}

\gksubpart{（1）}

\Needspace{5\baselineskip}
由数量积定义
\nopagebreak[4]
\[\overrightarrow{BA}\cdot\overrightarrow{BC}
=|\overrightarrow{BA}|\cdot|\overrightarrow{BC}|\cos B
=ca\cdot\frac13=2,\]
\Needspace{5\baselineskip}
得
\nopagebreak[4]
\[ac=6.\]
\Needspace{5\baselineskip}
由余弦定理
\nopagebreak[4]
\[b^2=a^2+c^2-2ac\cos B,\]
\Needspace{5\baselineskip}
即
\nopagebreak[4]
\[9=a^2+c^2-2\times6\times\frac13=a^2+c^2-4,\]
\Needspace{5\baselineskip}
故
\nopagebreak[4]
\[a^2+c^2=13.\]
\Needspace{5\baselineskip}
于是
\nopagebreak[4]
\[(a+c)^2=a^2+c^2+2ac=13+12=25,\]
\Needspace{5\baselineskip}
结合 $a+c>0$ 得
\nopagebreak[4]
\[a+c=5.\]
\Needspace{5\baselineskip}
所以 $a$，$c$ 是方程
\nopagebreak[4]
\[t^2-5t+6=0\]
\Needspace{5\baselineskip}
的两根，即 $t=2$ 或 $t=3$。又 $a>c$，故
\nopagebreak[4]
\[a=3,\qquad c=2.\]

\gksubpart{（2）}

\Needspace{5\baselineskip}
由余弦定理
\nopagebreak[4]
\[\cos C=\frac{a^2+b^2-c^2}{2ab}
=\frac{9+9-4}{2\times3\times3}
=\frac{14}{18}=\frac79.\]
\Needspace{5\baselineskip}
因为 $C\in(0,\pi)$，所以
\nopagebreak[4]
\[\sin C=\sqrt{1-\cos^2C}
=\sqrt{1-\frac{49}{81}}
=\frac{\sqrt{32}}9=\frac{4\sqrt2}9.\]
\Needspace{5\baselineskip}
又 $B\in(0,\pi)$，$\cos B=\dfrac13$，故
\nopagebreak[4]
\[\sin B=\sqrt{1-\frac19}=\frac{2\sqrt2}3.\]
\Needspace{5\baselineskip}
于是
\nopagebreak[4]
\[\begin{aligned}
\cos(B-C)
&=\cos B\cos C+\sin B\sin C\\
&=\frac13\cdot\frac79+\frac{2\sqrt2}3\cdot\frac{4\sqrt2}9\\
&=\frac7{27}+\frac{16}{27}=\frac{23}{27}.
\end{aligned}\]

\gkanswer{答案（1）：$a=3$，$c=2$。由 $ac\cos B=2$ 得 $ac=6$，再由余弦定理得 $a^2+c^2=13$，结合 $a>c$ 解出。}

\gkanswer{答案（2）：$\dfrac{23}{27}$。$\cos C=\dfrac79$，$\sin C=\dfrac{4\sqrt2}9$，$\sin B=\dfrac{2\sqrt2}3$，代入两角差余弦公式。}

\end{answer}
\begin{answer}
\gksolhead{18}{}

\textbf{题目}：

一家面包房根据以往某种面包的销售记录，绘制了日销售量的频率分布直方图，如图所示。将日销售量落入各组的频率视为概率，并假设每天的销售量相互独立。
\nopagebreak[4]
\begin{center}% ============================================================
% f2014_liaoning_science_q18_histogram.tex
% 2014 辽宁卷（理科）第 18 题：频率分布直方图
% 来源 PDF：普通高考/2014/2014辽宁理.pdf
% 原页：第 1 页（A3），印面 895，第 18 题题干右侧
% 组距 50，各组频率/组距依次为 0.003、0.005、0.006、0.004、0.002。
% ============================================================
\begin{tikzpicture}[>=Stealth, line width=0.7pt, font=\small,
  x=1.25cm, y=1cm, every node/.style={inner sep=1.5pt}]
  % 参考虚线（先画，随后被较高的矩形遮住，与原图一致）
  \foreach \h/\xend in {6/3, 5/2, 4/4, 2/5} {
    \draw[dashed, line width=0.35pt] (0,\h) -- (\xend,\h);
  }

  % 五个矩形
  \foreach \i/\h in {0/3, 1/5, 2/6, 3/4, 4/2} {
    \fill[white] (\i,0) rectangle (\i+1,\h);
    \draw (\i,\h) rectangle (\i+1,0);
  }

  % 坐标轴
  \draw[->] (-0.40,0) -- (6.55,0) node[right] {日销售量/个};
  \draw[->] (0,-0.25) -- (0,7.05);
  \node[align=center] at (-0.52,6.65) {频率\\[-2pt]组距};

  % 横轴刻度
  \foreach \i/\v in {0/0, 1/50, 2/100, 3/150, 4/200, 5/250} {
    \draw (\i,0.09) -- (\i,-0.09) node[below, font=\scriptsize] {\v};
  }

  % 纵轴刻度
  \foreach \h/\v in {2/0.002, 3/0.003, 4/0.004, 5/0.005, 6/0.006} {
    \draw (-0.09,\h) -- (0.09,\h);
    \node[left, font=\scriptsize] at (-0.11,\h) {\v};
  }
\end{tikzpicture}
\end{center}
\nopagebreak[4]
（1） 求在未来连续 3 天里，有连续 2 天的日销售量都不低于 100 个且另 1 天的日销售量低于 50 个的概率；

（2） 用 $X$ 表示在未来 3 天里日销售量不低于 100 个的天数，求随机变量 $X$ 的分布列，期望 $E(X)$ 及方差 $D(X)$。

\Needspace{4\baselineskip}
\textbf{解析}

\Needspace{5\baselineskip}
由频率分布直方图，组距为 $50$，各组的“频率/组距”依次为
\nopagebreak[4]
\[0.003,\quad0.005,\quad0.006,\quad0.004,\quad0.002.\]
\Needspace{5\baselineskip}
故各组的频率依次为
\nopagebreak[4]
\[\begin{aligned}
&0.003\times50=0.15,\\
&0.005\times50=0.25,\\
&0.006\times50=0.30,\\
&0.004\times50=0.20,\\
&0.002\times50=0.10.
\end{aligned}\]
\Needspace{5\baselineskip}
即日销售量不低于 $100$ 个的概率为
\nopagebreak[4]
\[p=0.30+0.20+0.10=0.6,\]
\Needspace{5\baselineskip}
低于 $50$ 个的概率为
\nopagebreak[4]
\[q=0.15.\]

\gksubpart{（1）}

\Needspace{5\baselineskip}
设一天日销售量不低于 $100$ 个为事件 $A$，低于 $50$ 个为事件 $B$，则
\nopagebreak[4]
\[P(A)=0.6,\qquad P(B)=0.15.\]

\Needspace{5\baselineskip}
连续 $3$ 天中，“有连续 $2$ 天 $A$ 发生且另 $1$ 天 $B$ 发生”的情形只有
\nopagebreak[4]
\[AAB,\qquad BAA\]
两种（$ABA$ 中的两个 $A$ 不相邻，不符合题意）。

\Needspace{5\baselineskip}
每种情形发生的概率为
\nopagebreak[4]
\[0.6\times0.6\times0.15=0.054.\]
\Needspace{5\baselineskip}
故所求概率为
\nopagebreak[4]
\[2\times0.054=0.108.\]

\gksubpart{（2）}

\Needspace{5\baselineskip}
每天日销售量不低于 $100$ 个的概率为 $0.6$，且各天相互独立，故
\nopagebreak[4]
\[X\sim B(3,0.6).\]
\Needspace{5\baselineskip}
分布列为
\nopagebreak[4]
\[\begin{aligned}
P(X=0)&=\mathrm C_3^0(0.6)^0(0.4)^3=0.064,\\
P(X=1)&=\mathrm C_3^1(0.6)^1(0.4)^2=0.288,\\
P(X=2)&=\mathrm C_3^2(0.6)^2(0.4)^1=0.432,\\
P(X=3)&=\mathrm C_3^3(0.6)^3(0.4)^0=0.216.
\end{aligned}\]
\Needspace{5\baselineskip}
所以
\nopagebreak[4]
\[E(X)=np=3\times0.6=1.8,\]
\nopagebreak[4]
\[D(X)=np(1-p)=3\times0.6\times0.4=0.72.\]

\gkanswer{答案（1）：$0.108$。日销售量不低于 $100$ 的概率为 $0.6$，低于 $50$ 的概率为 $0.15$，所求为 $2\times0.6^2\times0.15=0.108$。}

\gkanswer{答案（2）：$P(X=0)=0.064$，$P(X=1)=0.288$，$P(X=2)=0.432$，$P(X=3)=0.216$；$E(X)=1.8$，$D(X)=0.72$。$X\sim B(3,0.6)$，直接代入二项分布公式。}

\end{answer}
\begin{answer}
\gksolhead{19}{}

\textbf{题目}：

如图，$\triangle ABC$ 和 $\triangle BCD$ 所在平面互相垂直，且 $AB=BC=BD=2$，$\angle ABC=\angle DBC=120^\circ$，$E$，$F$ 分别为 $AC$，$DC$ 的中点。
\nopagebreak[4]
\begin{center}% ============================================================
% f2014_liaoning_science_q19_tetrahedron.tex
% 2014 辽宁卷（理科）第 19 题：两个平面垂直的四面体
% 来源 PDF：普通高考/2014/2014辽宁理.pdf
% 原页：第 2 页（A3），印面 896，第 19 题题干下方
% 图中 E、F 分别为 AC、DC 的中点；AB、BC、BD、BE 为被遮住的棱，
% 画虚线。
% ============================================================
\begin{tikzpicture}[>=Stealth, line width=0.7pt, font=\small, x=1cm, y=1cm,
  every node/.style={inner sep=1.5pt}]
  \coordinate (A) at (2.00,4.00);
  \coordinate (C) at (4.90,2.20);
  \coordinate (D) at (1.35,0.35);
  \coordinate (B) at (2.75,1.95);
  \coordinate (E) at (3.45,3.10);
  \coordinate (F) at (3.125,1.275);

  % 可见棱与截面线
  \draw (A) -- (C) -- (D) -- cycle;
  \draw (E) -- (F);

  % 被遮住的棱
  \draw[dashed] (A) -- (B);
  \draw[dashed] (B) -- (C);
  \draw[dashed] (B) -- (D);
  \draw[dashed] (B) -- (E);

  % 顶点与中点
  \fill (A) circle (1.1pt) node[above] {$A$};
  \fill (C) circle (1.1pt) node[right] {$C$};
  \fill (D) circle (1.1pt) node[below left] {$D$};
  \fill (B) circle (1.1pt) node[below right] {$B$};
  \fill (E) circle (1.1pt) node[above right] {$E$};
  \fill (F) circle (1.1pt) node[below right] {$F$};
\end{tikzpicture}
\end{center}
\nopagebreak[4]
（1） 求证：$EF\perp BC$；

（2） 求二面角 $E-BF-C$ 的正弦值。

\Needspace{4\baselineskip}
\textbf{解析}

建立空间直角坐标系。以 $B$ 为原点，$BC$ 所在直线为 $x$ 轴。因为平面 $ABC\perp$ 平面 $BCD$，可取平面 $ABC$ 为 $xOy$ 平面，平面 $BCD$ 为 $xOz$ 平面。

\Needspace{5\baselineskip}
由 $BC=2$ 得
\nopagebreak[4]
\[C(2,0,0).\]
\Needspace{5\baselineskip}
由 $AB=2$，$\angle ABC=120^\circ$，得
\nopagebreak[4]
\[A(2\cos120^\circ,2\sin120^\circ,0)=(-1,\sqrt3,0).\]
\Needspace{5\baselineskip}
由 $BD=2$，$\angle DBC=120^\circ$，且 $D$ 在 $xOz$ 平面内，得
\nopagebreak[4]
\[D(2\cos120^\circ,0,2\sin120^\circ)=(-1,0,\sqrt3).\]
\Needspace{5\baselineskip}
于是
\nopagebreak[4]
\[E\left(\frac12,\frac{\sqrt3}2,0\right),\qquad
F\left(\frac12,0,\frac{\sqrt3}2\right).\]

\gksubpart{（1）}

\Needspace{5\baselineskip}
计算
\nopagebreak[4]
\[\overrightarrow{EF}
=F-E=\left(0,-\frac{\sqrt3}2,\frac{\sqrt3}2\right),\]
\nopagebreak[4]
\[\overrightarrow{BC}=C-B=(2,0,0).\]
\Needspace{5\baselineskip}
所以
\nopagebreak[4]
\[\overrightarrow{EF}\cdot\overrightarrow{BC}
=0\cdot2+\left(-\frac{\sqrt3}2\right)\cdot0+\frac{\sqrt3}2\cdot0=0,\]
\Needspace{5\baselineskip}
故
\nopagebreak[4]
\[EF\perp BC.\]

\gksubpart{（2）}

\Needspace{5\baselineskip}
平面 $BEF$ 过原点 $B$，其法向量可取
\nopagebreak[4]
\[\boldsymbol n_1=\overrightarrow{BE}\times\overrightarrow{BF}.\]
\Needspace{5\baselineskip}
而
\nopagebreak[4]
\[\overrightarrow{BE}=\left(\frac12,\frac{\sqrt3}2,0\right),\qquad
\overrightarrow{BF}=\left(\frac12,0,\frac{\sqrt3}2\right).\]
\Needspace{5\baselineskip}
计算
\nopagebreak[4]
\[\begin{aligned}
\boldsymbol n_1
&=\begin{vmatrix}
\boldsymbol i&\boldsymbol j&\boldsymbol k\\
\frac12&\frac{\sqrt3}2&0\\
\frac12&0&\frac{\sqrt3}2
\end{vmatrix}\\
&=\left(\frac34,-\frac{\sqrt3}4,-\frac{\sqrt3}4\right)
\parallel(3,-\sqrt3,-\sqrt3).
\end{aligned}\]

\Needspace{5\baselineskip}
平面 $BCF$ 过原点 $B$，其法向量可取
\nopagebreak[4]
\[\boldsymbol n_2=\overrightarrow{BC}\times\overrightarrow{BF}.\]
\Needspace{5\baselineskip}
计算
\nopagebreak[4]
\[\begin{aligned}
\boldsymbol n_2
&=(2,0,0)\times\left(\frac12,0,\frac{\sqrt3}2\right)\\
&=(0,-\sqrt3,0)\parallel(0,1,0).
\end{aligned}\]

\Needspace{5\baselineskip}
设二面角 $E-BF-C$ 的大小为 $\theta$，则
\nopagebreak[4]
\[\begin{aligned}
|\cos\theta|
&=\frac{|\boldsymbol n_1\cdot\boldsymbol n_2|}
{|\boldsymbol n_1||\boldsymbol n_2|}\\
&=\frac{\sqrt3}{\sqrt{9+3+3}\cdot1}
=\frac{\sqrt3}{\sqrt{15}}
=\frac1{\sqrt5}.
\end{aligned}\]
\Needspace{5\baselineskip}
所以
\nopagebreak[4]
\[\sin\theta=\sqrt{1-\cos^2\theta}
=\sqrt{1-\frac15}
=\frac2{\sqrt5}
=\frac{2\sqrt5}5.\]

\gkanswer{答案（1）：见解析版}

\gkanswer{答案（2）：$\dfrac{2\sqrt5}5$。平面 $BEF$ 与平面 $BCF$ 的法向量分别为 $(3,-\sqrt3,-\sqrt3)$ 与 $(0,1,0)$，二面角余弦的绝对值为 $\dfrac1{\sqrt5}$，故正弦为 $\dfrac{2\sqrt5}5$。}

\end{answer}
\begin{answer}
\gksolhead{20}{}

\textbf{题目}：

圆 $x^2+y^2=4$ 的切线与 $x$ 轴正半轴，$y$ 轴正半轴围成一个三角形，当该三角形面积最小时，切点为 $P$（如图）。双曲线 $C_1:\dfrac{x^2}{a^2}-\dfrac{y^2}{b^2}=1$ 过点 $P$ 且离心率为 $\sqrt3$。
\nopagebreak[4]
\begin{center}% ============================================================
% f2014_liaoning_science_q20_circle_tangent.tex
% 2014 辽宁卷（理科）第 20 题：圆与第一象限内的切线
% 来源 PDF：普通高考/2014/2014辽宁理.pdf
% 原页：第 2 页（A3），印面 896，第 20 题题干下方
% 圆 x^2+y^2=4 的切线与两坐标轴正半轴围成三角形，切点为 P。
% 图中按面积最小的情形取 P(sqrt2,sqrt2)，切线 x+y=2sqrt2。
% ============================================================
\begin{tikzpicture}[>=Stealth, line width=0.7pt, font=\small, x=1cm, y=1cm,
  every node/.style={inner sep=1.5pt}]
  % 坐标轴
  \draw[->] (-3.1,0) -- (3.9,0) node[right] {$x$};
  \draw[->] (0,-3.1) -- (0,3.9) node[above] {$y$};
  \node[below left] at (0,0) {$O$};

  % 圆
  \draw (0,0) circle (2);

  % 切线与切点
  \draw (2.828,0) -- (0,2.828);
  \fill (1.414,1.414) circle (1.2pt) node[above right] {$P$};
\end{tikzpicture}
\end{center}
\nopagebreak[4]
（1） 求 $C_1$ 的方程；

（2） 椭圆 $C_2$ 过点 $P$ 且与 $C_1$ 有相同的焦点，直线 $l$ 过 $C_2$ 的右焦点且与 $C_2$ 交于 $A$，$B$ 两点，若以线段 $AB$ 为直径的圆过点 $P$，求 $l$ 的方程。

\Needspace{4\baselineskip}
\textbf{解析}

\gksubpart{（1）}

\Needspace{5\baselineskip}
设切点为 $P(x_0,y_0)$，$x_0>0$，$y_0>0$。圆 $x^2+y^2=4$ 在 $P$ 处的切线为
\nopagebreak[4]
\[x_0x+y_0y=4,\]
\Needspace{5\baselineskip}
即
\nopagebreak[4]
\[\frac{x}{4/x_0}+\frac{y}{4/y_0}=1.\]
\Needspace{5\baselineskip}
它与两坐标轴正半轴围成的三角形面积为
\nopagebreak[4]
\[S=\frac12\cdot\frac4{x_0}\cdot\frac4{y_0}
=\frac8{x_0y_0}.\]
\Needspace{5\baselineskip}
由
\nopagebreak[4]
\[x_0^2+y_0^2=4\geqslant2x_0y_0\]
\Needspace{5\baselineskip}
得
\nopagebreak[4]
\[x_0y_0\leqslant2,\]
\Needspace{5\baselineskip}
故
\nopagebreak[4]
\[S\geqslant\frac82=4,\]
\Needspace{5\baselineskip}
等号当且仅当
\nopagebreak[4]
\[x_0=y_0=\sqrt2.\]
\Needspace{5\baselineskip}
所以面积最小时
\nopagebreak[4]
\[P(\sqrt2,\sqrt2).\]

\Needspace{5\baselineskip}
双曲线 $C_1$ 的离心率 $e=\sqrt3$，故
\nopagebreak[4]
\[e^2=\frac{c^2}{a^2}=1+\frac{b^2}{a^2}=3,\]
\Needspace{5\baselineskip}
即
\nopagebreak[4]
\[b^2=2a^2.\]
\Needspace{5\baselineskip}
又 $P(\sqrt2,\sqrt2)$ 在 $C_1$ 上，
\nopagebreak[4]
\[\frac2{a^2}-\frac2{b^2}=1.\]
\Needspace{5\baselineskip}
代入 $b^2=2a^2$：
\nopagebreak[4]
\[\frac2{a^2}-\frac1{a^2}=1,\]
\Needspace{5\baselineskip}
故
\nopagebreak[4]
\[a^2=1,\qquad b^2=2.\]
\Needspace{5\baselineskip}
所以
\nopagebreak[4]
\[C_1:x^2-\frac{y^2}2=1.\]

\gksubpart{（2）}

\Needspace{5\baselineskip}
$C_1$ 的半焦距为
\nopagebreak[4]
\[c_1=\sqrt{a^2+b^2}=\sqrt3,\]
\Needspace{5\baselineskip}
故 $C_1$ 的焦点为 $(\pm\sqrt3,0)$。设椭圆
\nopagebreak[4]
\[C_2:\frac{x^2}{m}+\frac{y^2}{n}=1\quad(m>n>0).\]
\Needspace{5\baselineskip}
它与 $C_1$ 有相同焦点，故
\nopagebreak[4]
\[m-n=3.\]
\Needspace{5\baselineskip}
又 $P(\sqrt2,\sqrt2)$ 在 $C_2$ 上，
\nopagebreak[4]
\[\frac2m+\frac2n=1.\]
\Needspace{5\baselineskip}
由 $m=n+3$ 得
\nopagebreak[4]
\[\frac2{n+3}+\frac2n=1
\;\Longrightarrow\;
2n+2(n+3)=n(n+3),\]
\Needspace{5\baselineskip}
即
\nopagebreak[4]
\[n^2-n-6=0,\]
\Needspace{5\baselineskip}
解得 $n=3$ 或 $n=-2$。取 $n=3$，则 $m=6$。所以
\nopagebreak[4]
\[C_2:\frac{x^2}6+\frac{y^2}3=1,\]
右焦点为 $F(\sqrt3,0)$。

\Needspace{5\baselineskip}
设 $l:y=k(x-\sqrt3)$。代入 $x^2+2y^2=6$：
\nopagebreak[4]
\[(1+2k^2)x^2-4\sqrt3k^2x+6k^2-6=0.\]
\Needspace{5\baselineskip}
设 $A(x_1,y_1)$，$B(x_2,y_2)$，则
\nopagebreak[4]
\[x_1+x_2=\frac{4\sqrt3k^2}{1+2k^2},\qquad
x_1x_2=\frac{6k^2-6}{1+2k^2}.\]
\Needspace{5\baselineskip}
令
\nopagebreak[4]
\[u_i=x_i-\sqrt3,\]
\Needspace{5\baselineskip}
则
\nopagebreak[4]
\[u_1+u_2=x_1+x_2-2\sqrt3
=-\frac{2\sqrt3}{1+2k^2},\]
\nopagebreak[4]
\[u_1u_2=x_1x_2-\sqrt3(x_1+x_2)+3
=-\frac3{1+2k^2}.\]

\Needspace{5\baselineskip}
以 $AB$ 为直径的圆过 $P$，等价于
\nopagebreak[4]
\[\overrightarrow{PA}\cdot\overrightarrow{PB}=0.\]
\Needspace{5\baselineskip}
而
\nopagebreak[4]
\[\overrightarrow{PA}=(x_1-\sqrt2,y_1-\sqrt2),\qquad
\overrightarrow{PB}=(x_2-\sqrt2,y_2-\sqrt2),\]
\Needspace{5\baselineskip}
且 $y_i-\sqrt2=ku_i-\sqrt2$，$x_i-\sqrt2=u_i+(\sqrt3-\sqrt2)$。于是
\nopagebreak[4]
\[\begin{aligned}
\overrightarrow{PA}\cdot\overrightarrow{PB}
&=(1+k^2)u_1u_2
+\bigl(\sqrt3-\sqrt2-\sqrt2k\bigr)(u_1+u_2)\\
&\qquad+(\sqrt3-\sqrt2)^2+2.
\end{aligned}\]
\Needspace{5\baselineskip}
代入 $u_1+u_2$、$u_1u_2$，两边乘以 $1+2k^2$：
\nopagebreak[4]
\[\begin{aligned}
&-3(1+k^2)-2\sqrt3\bigl(\sqrt3-\sqrt2-\sqrt2k\bigr)\\
&\qquad+(7-2\sqrt6)(1+2k^2)=0.
\end{aligned}\]
\Needspace{5\baselineskip}
整理得
\nopagebreak[4]
\[(11-4\sqrt6)k^2+2\sqrt6k-2=0.\]
\Needspace{5\baselineskip}
其判别式
\nopagebreak[4]
\[\begin{aligned}
\Delta
&=(2\sqrt6)^2+8(11-4\sqrt6)\\
&=24+88-32\sqrt6\\
&=112-32\sqrt6\\
&=16(7-2\sqrt6)=16(\sqrt6-1)^2.
\end{aligned}\]
\Needspace{5\baselineskip}
故
\nopagebreak[4]
\[\begin{aligned}
k&=\frac{-2\sqrt6\pm4(\sqrt6-1)}{2(11-4\sqrt6)}\\
&=\frac{3\sqrt6+2}{25}\quad\text{\text{或}}\quad
k=-(2+\sqrt6).
\end{aligned}\]

\Needspace{5\baselineskip}
再检查竖直线 $x=\sqrt3$。此时 $A$、$B$ 的纵坐标为 $\pm\sqrt{\dfrac32}$，于是
\nopagebreak[4]
\[\overrightarrow{PA}\cdot\overrightarrow{PB}
=(\sqrt3-\sqrt2)^2+\frac12
=\frac{11}2-2\sqrt6\neq0,\]
故竖直线不是解。

\Needspace{5\baselineskip}
所以
\nopagebreak[4]
\[l:y=\frac{3\sqrt6+2}{25}(x-\sqrt3)\]
\Needspace{5\baselineskip}
或
\nopagebreak[4]
\[l:y=-(2+\sqrt6)(x-\sqrt3).\]

\gkanswer{答案（1）：$x^2-\dfrac{y^2}2=1$。由 $e=\sqrt3$ 得 $b^2=2a^2$，代入 $P(\sqrt2,\sqrt2)$ 得 $a^2=1$。}

\gkanswer{答案（2）：$y=\dfrac{3\sqrt6+2}{25}(x-\sqrt3)$ 或 $y=-(2+\sqrt6)(x-\sqrt3)$。设 $l:y=k(x-\sqrt3)$，由 $\overrightarrow{PA}\cdot\overrightarrow{PB}=0$ 得 $(11-4\sqrt6)k^2+2\sqrt6k-2=0$，解出 $k$。}

\end{answer}
\begin{answer}
\gksolhead{21}{}

\textbf{题目}：

已知函数 $f(x)=(\cos x-x)(\pi+2x)-\dfrac83(\sin x+1)$，$g(x)=3(x-\pi)\cos x-4(1+\sin x)\ln\left(3-\dfrac{2x}{\pi}\right)$。

（1） 证明：存在唯一 $x_0\in\left(0,\dfrac\pi2\right)$，使 $f(x_0)=0$；

（2） 证明：存在唯一 $x_1\in\left(\dfrac\pi2,\pi\right)$，使 $g(x_1)=0$，且对（1）中的 $x_0$，有 $x_0+x_1<\pi$。

\Needspace{4\baselineskip}
\textbf{解析}

\gksubpart{（1）}

\Needspace{5\baselineskip}
求导：
\nopagebreak[4]
\[\begin{aligned}
f'(x)
&=(-\sin x-1)(\pi+2x)+2(\cos x-x)-\frac83\cos x\\
&=-(1+\sin x)(\pi+2x)-2x-\frac23\cos x.
\end{aligned}\]
\Needspace{5\baselineskip}
当 $0<x<\dfrac\pi2$ 时，
\nopagebreak[4]
\[1+\sin x>0,\qquad \pi+2x>0,\qquad x>0,\qquad \cos x>0,\]
\Needspace{5\baselineskip}
故
\nopagebreak[4]
\[f'(x)<0.\]
所以 $f$ 在 $\left(0,\dfrac\pi2\right)$ 上严格单调递减。

\Needspace{5\baselineskip}
又
\nopagebreak[4]
\[f(0)=\pi-\frac83>0,\]
\nopagebreak[4]
\[f\left(\frac\pi2\right)
=\left(0-\frac\pi2\right)(2\pi)-\frac83(1+1)
=-\pi^2-\frac{16}3<0.\]
由零点存在定理，存在 $x_0\in\left(0,\dfrac\pi2\right)$ 使 $f(x_0)=0$；由严格单调性，这样的 $x_0$ 唯一。

\gksubpart{（2）}

\Needspace{5\baselineskip}
先证存在性。计算
\nopagebreak[4]
\[g\left(\frac\pi2\right)
=3\left(-\frac\pi2\right)\cdot0-4(2)\ln2
=-8\ln2<0.\]
\Needspace{5\baselineskip}
又
\nopagebreak[4]
\[g(\pi)=3\cdot0\cdot(-1)-4\cdot1\cdot\ln1=0,\]
\Needspace{5\baselineskip}
且
\nopagebreak[4]
\[\begin{aligned}
g'(\pi)
&=3\cos\pi-3(\pi-\pi)\sin\pi-4\cos\pi\ln1
+\frac{8(1+\sin\pi)}{\pi(3-2)}\\
&=-3+\frac8\pi<0.
\end{aligned}\]
因为 $g(\pi)=0$ 且 $g'(\pi)<0$，所以在 $\pi$ 左侧充分近处有 $g(x)>0$。结合 $g\left(\dfrac\pi2\right)<0$，由零点存在定理，$g$ 在 $\left(\dfrac\pi2,\pi\right)$ 内有零点。

\Needspace{5\baselineskip}
再证唯一性。令
\nopagebreak[4]
\[u=\pi-x,\qquad x\in\left(\frac\pi2,\pi\right)
\;\Longleftrightarrow\;
u\in\left(0,\frac\pi2\right),\]
\Needspace{5\baselineskip}
并令
\nopagebreak[4]
\[G(u)=g(\pi-u)
=3u\cos u-4(1+\sin u)\ln\left(1+\frac{2u}{\pi}\right).\]
则 $g$ 在 $\left(\dfrac\pi2,\pi\right)$ 内的零点与 $G$ 在 $\left(0,\dfrac\pi2\right)$ 内的零点一一对应。

\Needspace{5\baselineskip}
求导：
\nopagebreak[4]
\[\begin{aligned}
G'(u)
&=3\cos u-3u\sin u-4\cos u\ln\left(1+\frac{2u}{\pi}\right)\\
&\qquad-\frac{8(1+\sin u)}{\pi+2u},\\
G''(u)
&=-\sin u\left[6-4\ln\left(1+\frac{2u}{\pi}\right)\right]-3u\cos u\\
&\qquad-\frac{16\cos u}{\pi+2u}
+\frac{16(1+\sin u)}{(\pi+2u)^2}.
\end{aligned}\]
\Needspace{5\baselineskip}
由于
\nopagebreak[4]
\[\ln\left(1+\frac{2u}{\pi}\right)\leqslant\ln2,\]
\Needspace{5\baselineskip}
故
\nopagebreak[4]
\[6-4\ln\left(1+\frac{2u}{\pi}\right)\geqslant6-4\ln2>3.\]
\Needspace{5\baselineskip}
对 $0<u\leqslant\dfrac12$，利用 $\cos u\geqslant\cos\dfrac12$、$\pi+2u\leqslant\pi+1$；对 $\dfrac12<u\leqslant1$，利用 $\sin u\geqslant\sin\dfrac12$、$\cos u\geqslant\cos1$；对 $1<u<\dfrac\pi2$，利用 $\sin u\geqslant\sin1$、$\pi+2u\geqslant\pi+2$。再注意 $1+\sin u\leqslant2$，分段放缩均得
\nopagebreak[4]
\[G''(u)<0.\]
\Needspace{5\baselineskip}
（具体地，
\nopagebreak[4]
\[\begin{aligned}
0<u\leqslant\tfrac12:\quad
-G''(u)&\geqslant\frac{16\cos\frac12}{\pi+1}-\frac{32}{\pi^2}>0,\\
\tfrac12<u\leqslant1:\quad
-G''(u)&\geqslant3\sin\frac12+\frac{16\cos1}{\pi+2}-\frac{32}{(\pi+1)^2}>0,\\
1<u<\tfrac\pi2:\quad
-G''(u)&\geqslant3\sin1-\frac{32}{(\pi+2)^2}>0.
\end{aligned}\]
）
所以 $G$ 在 $\left(0,\dfrac\pi2\right)$ 上严格上凸。

\Needspace{5\baselineskip}
又
\nopagebreak[4]
\[G(0)=0,\qquad
G\left(\frac\pi2\right)=-8\ln2<0,\]
\nopagebreak[4]
\[G'(0)=3-\frac8\pi>0.\]
上凸函数 $G$ 从 $G(0)=0$ 出发先增后减，且 $G\left(\dfrac\pi2\right)<0$，故它在 $\left(0,\dfrac\pi2\right)$ 内恰有一个零点。于是 $g$ 在 $\left(\dfrac\pi2,\pi\right)$ 内恰有一个零点，记为 $x_1$。

最后证明 $x_0+x_1<\pi$。

\Needspace{5\baselineskip}
令
\nopagebreak[4]
\[t=x_0,\qquad s=\frac{2t}{\pi}\in(0,1),\qquad L=\ln(1+s).\]
\Needspace{5\baselineskip}
由 $f(t)=0$ 得
\nopagebreak[4]
\[(\cos t-t)(\pi+2t)=\frac83(1+\sin t).\tag{1}\]
\Needspace{5\baselineskip}
于是
\nopagebreak[4]
\[\begin{aligned}
g(\pi-t)
&=3t\cos t-4(1+\sin t)L\\
&=3t\cos t-\frac32(\cos t-t)(\pi+2t)L\\
&=3\left[t\cos t-\frac12(\cos t-t)(\pi+2t)L\right].
\end{aligned}\]
\Needspace{5\baselineskip}
因此 $g(\pi-t)>0$ 等价于
\nopagebreak[4]
\[2t\cos t>(\cos t-t)(\pi+2t)L.\]
\Needspace{5\baselineskip}
两边除以 $\pi$，并利用 $t=\dfrac{\pi s}{2}$、$\pi+2t=\pi(1+s)$，得
\nopagebreak[4]
\[s\cos t>(1+s)L\left(\cos t-\frac{\pi s}{2}\right),\]
\Needspace{5\baselineskip}
即
\nopagebreak[4]
\[(1+s)Lt>\cos t\bigl[(1+s)L-s\bigr].\tag{2}\]

\Needspace{5\baselineskip}
考察函数
\nopagebreak[4]
\[\varphi(s)=(1+s)\ln(1+s)-s-\frac{s^2}{2}.\]
\Needspace{5\baselineskip}
则 $\varphi(0)=0$，且
\nopagebreak[4]
\[\varphi'(s)=\ln(1+s)-s<0\quad(s>0),\]
\Needspace{5\baselineskip}
故
\nopagebreak[4]
\[0<(1+s)L-s<\frac{s^2}{2}.\]
\Needspace{5\baselineskip}
又由
\nopagebreak[4]
\[\psi(s)=(1+s)\ln(1+s)-s,\qquad
\psi(0)=0,\qquad \psi'(s)=\ln(1+s)>0,\]
\Needspace{5\baselineskip}
得
\nopagebreak[4]
\[(1+s)L>s.\]
\Needspace{5\baselineskip}
于是（2）式右端
\nopagebreak[4]
\[\cos t\bigl[(1+s)L-s\bigr]<\frac{s^2}{2},\]
\Needspace{5\baselineskip}
而左端
\nopagebreak[4]
\[(1+s)Lt>st=\frac{\pi s^2}{2}>\frac{s^2}{2}.\]
\Needspace{5\baselineskip}
故（2）式成立，从而
\nopagebreak[4]
\[g(\pi-x_0)>0.\]

\Needspace{5\baselineskip}
因为 $g$ 在 $\left(\dfrac\pi2,\pi\right)$ 内只有一个零点 $x_1$，且 $g\left(\dfrac\pi2\right)<0$、$g(\pi-x_0)>0$，所以
\nopagebreak[4]
\[\pi-x_0>x_1,\]
\Needspace{5\baselineskip}
即
\nopagebreak[4]
\[x_0+x_1<\pi.\]

\gkanswer{答案（1）：见解析版}

\gkanswer{答案（2）：见解析版}

\end{answer}
\begin{answer}
\gksolhead{22}{}

\textbf{题目}：

【A】如图，$EP$ 交圆于 $E$，$C$ 两点，$PD$ 切圆于 $D$，$G$ 为 $CE$ 上一点且 $PG=PD$，连接 $DG$ 并延长交圆于点 $A$，作弦 $AB$ 垂直 $EP$，垂足为 $F$。
\nopagebreak[4]
\begin{center}% ============================================================
% f2014_liaoning_science_q22_circle.tex
% 2014 辽宁卷（理科）第 22 题：圆的切线与割线
% 来源 PDF：普通高考/2014/2014辽宁理.pdf
% 原页：第 2 页（A3），印面 896，第 22 题题干下方
% EP 交圆于 E、C，PD 切圆于 D，G 在 CE 上且 PG=PD，
% DG 延长交圆于 A，AB 垂直 EP 于 F。
% 取单位圆：E(-0.866,-0.5)、C(0.866,-0.5)、A(0,-1)、B(0,1)、
% D(0.574,0.819)、P(2.456,-0.5)，此时 PD 与圆相切、PG=PD。
% ============================================================
\begin{tikzpicture}[>=Stealth, line width=0.7pt, font=\small,
  x=1.5cm, y=1.5cm, every node/.style={inner sep=1.5pt}]
  \coordinate (E) at (-0.866,-0.5);
  \coordinate (C) at (0.866,-0.5);
  \coordinate (A) at (0,-1);
  \coordinate (B) at (0,1);
  \coordinate (F) at (0,-0.5);
  \coordinate (G) at (0.158,-0.5);
  \coordinate (D) at (0.574,0.819);
  \coordinate (P) at (2.456,-0.5);

  % 圆
  \draw (0,0) circle (1);

  % 割线 EP、直径 AB、弦 BD、切线 PD、弦 ED、直线 DGA、弦 AC
  \draw (E) -- (P);
  \draw (A) -- (B);
  \draw (B) -- (D);
  \draw (D) -- (P);
  \draw (E) -- (D);
  \draw (D) -- (A);
  \draw (A) -- (C);

  % 点与标签
  \fill (E) circle (1.1pt) node[left] {$E$};
  \fill (C) circle (1.1pt) node[below] {$C$};
  \fill (A) circle (1.1pt) node[below] {$A$};
  \fill (B) circle (1.1pt) node[above] {$B$};
  \fill (F) circle (1.0pt) node[above left, xshift=-2pt] {$F$};
  \fill (G) circle (1.0pt);
  \node at (0.23,-0.66) {$G$};
  \fill (D) circle (1.1pt) node[above right] {$D$};
  \fill (P) circle (1.1pt) node[right] {$P$};
\end{tikzpicture}
\end{center}
\nopagebreak[4]
（1） 求证：$AB$ 为圆的直径；

（2） 若 $AC=BD$，求证：$AB=ED$。

\Needspace{4\baselineskip}
\textbf{解析}

\gksubpart{（1）}

\Needspace{5\baselineskip}
因为 $PD$ 与圆相切于 $D$，由弦切角定理，切线 $DP$ 与弦 $DA$ 所成的、包含圆心的那个角等于 $\angle DCA$；而 $\angle PDA$ 与它互补，故
\nopagebreak[4]
\[\angle PDA+\angle DCA=180^\circ.\]
\Needspace{5\baselineskip}
又 $PD=PG$，故
\nopagebreak[4]
\[\angle PDG=\angle PGD.\]
\Needspace{5\baselineskip}
由于 $G$ 在 $DA$ 上，$\angle PDG=\angle PDA$，所以
\nopagebreak[4]
\[\angle PGD+\angle DCA=180^\circ.\]
\Needspace{5\baselineskip}
而 $P$，$G$，$C$ 共线，$D$，$G$，$A$ 共线且 $G$ 在 $D$ 与 $A$ 之间，故
\nopagebreak[4]
\[\angle PGD+\angle AGC=180^\circ.\]
\Needspace{5\baselineskip}
因此
\nopagebreak[4]
\[\angle AGC=\angle DCA.\]

\Needspace{5\baselineskip}
在 $\triangle GCD$ 中，$\angle AGC$ 是顶点 $G$ 处的外角，故
\nopagebreak[4]
\[\angle AGC=\angle GCD+\angle GDC.\]
\Needspace{5\baselineskip}
又 $CG$ 在 $\angle DCA$ 内部，
\nopagebreak[4]
\[\angle DCA=\angle GCD+\angle GCA.\]
\Needspace{5\baselineskip}
比较两式得
\nopagebreak[4]
\[\angle GCA=\angle GDC.\]
\Needspace{5\baselineskip}
因为 $G$，$D$，$A$ 共线，
\nopagebreak[4]
\[\angle GDC=\angle ADC.\]
\Needspace{5\baselineskip}
因为 $A$，$E$，$C$，$D$ 四点共圆，
\nopagebreak[4]
\[\angle ADC=\angle AEC.\]
\Needspace{5\baselineskip}
因为 $G$，$C$，$E$ 共线，
\nopagebreak[4]
\[\angle GCA=\angle ACE.\]
\Needspace{5\baselineskip}
所以
\nopagebreak[4]
\[\angle ACE=\angle AEC,\]
\Needspace{5\baselineskip}
从而
\nopagebreak[4]
\[AE=AC.\]

\Needspace{5\baselineskip}
在等腰三角形 $AEC$ 中，$AF\perp EC$，故 $AF$ 也是底边 $EC$ 的中线，即
\nopagebreak[4]
\[FE=FC.\]
于是 $AB$ 垂直平分弦 $EC$。弦的垂直平分线过圆心，故 $AB$ 过圆心；又 $A$，$B$ 在圆上，所以 $AB$ 为圆的直径。

\gksubpart{（2）}

\Needspace{5\baselineskip}
由（1）知 $AB$ 为直径，故
\nopagebreak[4]
\[\angle ACB=90^\circ,\qquad \angle ADB=90^\circ.\]
\Needspace{5\baselineskip}
在 $\mathrm{Rt}\triangle ACB$ 与 $\mathrm{Rt}\triangle BDA$ 中，
\nopagebreak[4]
\[AB=BA,\qquad AC=BD,\]
\Needspace{5\baselineskip}
由斜边、直角边对应相等得
\nopagebreak[4]
\[\mathrm{Rt}\triangle ACB\cong\mathrm{Rt}\triangle BDA.\]
\Needspace{5\baselineskip}
于是
\nopagebreak[4]
\[\angle ABC=\angle BAD.\]

\Needspace{5\baselineskip}
又由（1）知 $AB$ 垂直平分 $EC$，故 $AE=AC$，且 $AB$ 平分 $\angle EAC$，即
\nopagebreak[4]
\[\angle EAB=\angle BAC.\]
\Needspace{5\baselineskip}
因为 $AB$ 在 $\angle EAD$ 内部，
\nopagebreak[4]
\[\begin{aligned}
\angle EAD
&=\angle EAB+\angle BAD\\
&=\angle BAC+\angle ABC\\
&=180^\circ-\angle ACB\\
&=90^\circ.
\end{aligned}\]
\Needspace{5\baselineskip}
所以 $\angle EAD=90^\circ$，由圆周角定理的逆定理，$ED$ 为圆的直径。于是
\nopagebreak[4]
\[ED=AB.\]

\gkanswer{答案（1）：见解析版}

\gkanswer{答案（2）：见解析版}

\end{answer}
\begin{answer}
\gksolhead{23}{}

\textbf{题目}：

【B】将圆 $x^2+y^2=1$ 上每一点的横坐标保持不变，纵坐标变为原来的 $2$ 倍，得曲线 $C$。

（1） 写出 $C$ 的参数方程；

（2） 设直线 $l:2x+y-2=0$ 与 $C$ 的交点为 $P_1$，$P_2$，以坐标原点为极点，$x$ 轴正半轴为极轴建立极坐标系，求过线段 $P_1P_2$ 的中点且与 $l$ 垂直的直线的极坐标方程。

\Needspace{4\baselineskip}
\textbf{解析}

\gksubpart{（1）}

\Needspace{5\baselineskip}
圆 $x^2+y^2=1$ 上点的坐标可写为
\nopagebreak[4]
\[x=\cos\theta,\qquad y=\sin\theta\quad(\theta\in\mathbf R).\]
\Needspace{5\baselineskip}
纵坐标变为原来的 $2$ 倍、横坐标不变，得
\nopagebreak[4]
\[\begin{cases}
x=\cos\theta,\\
y=2\sin\theta,
\end{cases}\qquad\theta\in\mathbf R.\]
\Needspace{5\baselineskip}
消去 $\theta$ 得普通方程
\nopagebreak[4]
\[x^2+\frac{y^2}4=1.\]

\gksubpart{（2）}

\Needspace{5\baselineskip}
将 $y=2-2x$ 代入 $x^2+\dfrac{y^2}4=1$：
\nopagebreak[4]
\[x^2+\frac{(2-2x)^2}{4}=1,\]
\Needspace{5\baselineskip}
即
\nopagebreak[4]
\[x^2+(1-x)^2=1,\]
\Needspace{5\baselineskip}
化简得
\nopagebreak[4]
\[2x^2-2x=0,\]
\Needspace{5\baselineskip}
故 $x=0$ 或 $x=1$。于是两个交点为
\nopagebreak[4]
\[P_1(1,0),\qquad P_2(0,2).\]
\Needspace{5\baselineskip}
线段 $P_1P_2$ 的中点为
\nopagebreak[4]
\[M\left(\frac12,1\right).\]

\Needspace{5\baselineskip}
直线 $l:2x+y-2=0$ 的斜率为 $-2$，故与 $l$ 垂直的直线斜率为 $\dfrac12$，其方程为
\nopagebreak[4]
\[y-1=\frac12\left(x-\frac12\right),\]
\Needspace{5\baselineskip}
即
\nopagebreak[4]
\[2x-4y+3=0.\]
\Needspace{5\baselineskip}
代入极坐标变换
\nopagebreak[4]
\[x=\rho\cos\theta,\qquad y=\rho\sin\theta,\]
\Needspace{5\baselineskip}
得
\nopagebreak[4]
\[2\rho\cos\theta-4\rho\sin\theta+3=0,\]
\Needspace{5\baselineskip}
故
\nopagebreak[4]
\[\rho=\frac3{4\sin\theta-2\cos\theta}.\]

\gkanswer{答案（1）：$\begin{cases}x=\cos\theta,\\y=2\sin\theta\end{cases}$（$\theta$ 为参数）。由 $x^2+\dfrac{y^2}4=1$ 直接写出。}

\gkanswer{答案（2）：$\rho=\dfrac3{4\sin\theta-2\cos\theta}$。中点 $M\left(\dfrac12,1\right)$，所求直线 $2x-4y+3=0$，代入 $x=\rho\cos\theta$，$y=\rho\sin\theta$。}

\end{answer}
\begin{answer}
\gksolhead{24}{}

\textbf{题目}：

【C】设函数 $f(x)=2|x-1|+x-1$，$g(x)=16x^2-8x+1$，记 $f(x)\leqslant 1$ 的解集为 $M$，$g(x)\leqslant 4$ 的解集为 $N$。

（1） 求 $M$；

（2） 当 $x\in M\cap N$ 时，证明：$x^2f(x)+x[f(x)]^2\leqslant\dfrac14$。

\Needspace{4\baselineskip}
\textbf{解析}

\gksubpart{（1）}

\Needspace{5\baselineskip}
当 $x\geqslant1$ 时，
\nopagebreak[4]
\[f(x)=2(x-1)+x-1=3x-3.\]
\Needspace{5\baselineskip}
由 $f(x)\leqslant1$ 得
\nopagebreak[4]
\[3x-3\leqslant1,\qquad x\leqslant\frac43.\]
\Needspace{5\baselineskip}
结合 $x\geqslant1$，得
\nopagebreak[4]
\[1\leqslant x\leqslant\frac43.\]

\Needspace{5\baselineskip}
当 $x<1$ 时，
\nopagebreak[4]
\[f(x)=2(1-x)+x-1=1-x.\]
\Needspace{5\baselineskip}
由 $f(x)\leqslant1$ 得
\nopagebreak[4]
\[1-x\leqslant1,\qquad x\geqslant0.\]
\Needspace{5\baselineskip}
结合 $x<1$，得
\nopagebreak[4]
\[0\leqslant x<1.\]

\Needspace{5\baselineskip}
综上
\nopagebreak[4]
\[M=\left[0,\frac43\right].\]

\gksubpart{（2）}

\Needspace{5\baselineskip}
由
\nopagebreak[4]
\[g(x)=16x^2-8x+1=(4x-1)^2\leqslant4,\]
\Needspace{5\baselineskip}
得
\nopagebreak[4]
\[|4x-1|\leqslant2,\]
\Needspace{5\baselineskip}
即
\nopagebreak[4]
\[-2\leqslant4x-1\leqslant2,\]
\Needspace{5\baselineskip}
解得
\nopagebreak[4]
\[-\frac14\leqslant x\leqslant\frac34.\]
\Needspace{5\baselineskip}
故
\nopagebreak[4]
\[N=\left[-\frac14,\frac34\right],\]
\Needspace{5\baselineskip}
从而
\nopagebreak[4]
\[M\cap N=\left[0,\frac34\right].\]

\Needspace{5\baselineskip}
当 $x\in M\cap N$ 时，$0\leqslant x\leqslant\dfrac34<1$，故
\nopagebreak[4]
\[f(x)=1-x.\]
\Needspace{5\baselineskip}
于是
\nopagebreak[4]
\[\begin{aligned}
x^2f(x)+x[f(x)]^2
&=x^2(1-x)+x(1-x)^2\\
&=x(1-x)\bigl[x+(1-x)\bigr]\\
&=x(1-x).
\end{aligned}\]
\Needspace{5\baselineskip}
由均值不等式，
\nopagebreak[4]
\[x(1-x)\leqslant\left(\frac{x+(1-x)}2\right)^2=\frac14,\]
\Needspace{5\baselineskip}
等号在 $x=\dfrac12$ 时成立。故
\nopagebreak[4]
\[x^2f(x)+x[f(x)]^2\leqslant\frac14.\]

\gkanswer{答案（1）：$M=\left[0,\dfrac43\right]$。分 $x\geqslant1$ 与 $x<1$ 去绝对值求解。}

\gkanswer{答案（2）：见解析版}

\end{answer}